\section{AMC 10 10B 2025}
        \begin{enumerate}
        	\item (\textit{AMC 10 10B 2025}) The instructions on a $350$-gram bag of coffee beans say that proper brewing of a large mug of pour-over coffee requires $20$ grams of coffee beans. What is the greatest number of properly brewed large mugs of coffee that can be made from the coffee beans in that bag?


 $\textbf{(A) } 16 \qquad\textbf{(B) } 17 \qquad\textbf{(C) } 18 \qquad\textbf{(D) } 19 \qquad\textbf{(E) } 20$


	\item (\textit{AMC 10 10B 2025}) Jerry wrote down the ones digit of each of the first $2025$ positive squares: $1,4,9,6,5,6,\dots$. What is the sum of all the numbers Jerry wrote down?


 $\textbf{(A)}~9025 \qquad \textbf{(B)}~9070 \qquad \textbf{(C)}~9090 \qquad \textbf{(D)}~9115 \qquad \textbf{(E)}~9160$


	\item (\textit{AMC 10 10B 2025}) A Pascal-like triangle has $10$ as the top row and $10$ followed by $1$ as the second row. In each subsequent row the first number is $10$, the last number is $1$, and, as in the standard Pascal's Triangle, each other number in the row is the sum of the two numbers directly above it. The first four rows are shown below:


 $\begin{tabular}{ccccccc} &&&10&&&\\ &&10&&1&& \\ &10&&11&&1& \\ 10&&21&&12&&1 \\ \end{tabular}$


 What is the sum of the digits of the sum of the numbers in the $11$th row? 


 $\textbf{(A) } 11 \qquad\textbf{(B) } 13 \qquad\textbf{(C) } 14 \qquad\textbf{(D) } 16 \qquad\textbf{(E) } 17$


	\item (\textit{AMC 10 10B 2025}) The value of the two-digit number $\underline{a}~\underline{b}$ in base seven equals the value of the two-digit number $\underline{b}~\underline{a}$ in base nine. What is $a+b?$
$\textbf{(A)}~7\qquad\textbf{(B)}~9\qquad\textbf{(C)}~10\qquad\textbf{(D)}~11\qquad\textbf{(E)}~14$


	\item (\textit{AMC 10 10B 2025}) In $\triangle ABC$, $AB=10$, $AC=18$, and $\angle B=130^\circ.$ Let $O$ be the center of the circle containing $A$, $B$, and $C.$ What is the degree measure of $\angle CAO$?


 $\textbf{(A)}~20 \qquad \textbf{(B)}~30 \qquad \textbf{(C)}~40 \qquad \textbf{(D)}~50 \qquad \textbf{(E)}~60$


	\item (\textit{AMC 10 10B 2025}) The line $y = \frac{1}{3}x + 1$ divides the square region defined by $0 \leq x \leq 2$ and $0 \leq y \leq 2$ into an upper and lower region. The line $x=a$ divides the lower region into two regions of equal area. Then $a$ can be written as $\sqrt{s} - t$, where $s$ and $t$ are positive integers. What is $s+t$?


 $\textbf{(A) } 18 \qquad\textbf{(B) } 19 \qquad\textbf{(C) }20\qquad\textbf{(D) } 21 \qquad\textbf{(E) } 22$


	\item (\textit{AMC 10 10B 2025}) Frances stands $15$ meters directly south of a locked gate in a fence that runs east-west. Immediately behind the fence is a box of chocolates, located $x$ meters east of the locked gate. An unlocked gate lies $9$ meters east of the box, and another unlocked gate lies $8$ meters west of the locked gate. Frances can reach the box by walking toward an unlocked gate, passing through it, and walking toward the box. It happens that the total distance Frances would travel is the same via either unlocked gate. What is the value of $x$?


 $\textbf{(A)}~3\frac27 \qquad \textbf{(B)}~3\frac37 \qquad \textbf{(C)}~3\frac47 \qquad \textbf{(D)}~3\frac57 \qquad \textbf{(E)}~3\frac67$


	\item (\textit{AMC 10 10B 2025}) Emmy says to Max, "I ordered $36$ math club sweatshirts today." Max asks, "How much did each shirt cost?" Emmy responds, "I'll give you a hint. The total cost was $\$ \underline A~\underline B~\underline B.\underline B~\underline A$, where $A$ and $B$ are digits and $A \neq 0$." After a pause, Max says, "That was a good price." What is $A + B$?


 $\textbf{(A)}~7 \qquad \textbf{(B)}~8 \qquad \textbf{(C)}~11 \qquad \textbf{(D)}~14 \qquad \textbf{(E)}~15$


	\item (\textit{AMC 10 10B 2025}) How many ordered triples of integers $(x, y, z)$ satisfy the following system of inequalities?


 
 \begin{align*} -x-y-z&\le -2\\ -x+y+z&\le 2\\ x-y+z&\le 2\\ x+y-z&\le 2 \end{align*}

$\textbf{(A)}\ 4 \qquad \textbf{(B)}\ 8 \qquad \textbf{(C)}\ 11 \qquad \textbf{(D)}\ 15 \qquad \textbf{(E)}\ 17 \qquad$


	\item (\textit{AMC 10 10B 2025}) Let $f(n)=n^3-5n^2+2n+8$ and $g(n)=n^3-6n^2+5n+12.$ What is the sum of all integers $n$ such that $\frac{f(n)}{g(n)}$ is an integer?


 $\textbf{(A)}~2\qquad\textbf{(B)}~3\qquad\textbf{(C)}~4\qquad\textbf{(D)}~5\qquad\textbf{(E)}~6$


	\item (\textit{AMC 10 10B 2025}) On Monday, $6$ students went to the tutoring center at the same time, and each one was randomly assigned to one of the $6$ tutors on duty. On Tuesday, the same $6$ students showed up, the same $6$ tutors were on duty, and the students were again randomly assigned to the tutors. What is the probability that exactly $2$ students met with the same tutor both Monday and Tuesday?


 $\textbf{(A) } \frac{1}{16} \qquad\textbf{(B) } \frac{3}{16} \qquad\textbf{(C) }\frac{1}{4}\qquad\textbf{(D) } \frac{3}{8} \qquad\textbf{(E) } \frac{1}{2}$


	\item (\textit{AMC 10 10B 2025}) The figure below shows an equilateral triangle, a rhombus with a $60^\circ$ angle, and a regular hexagon, each of them containing some mutually tangent congruent disks. Let $T, R,$ and $H,$ respectively, denote the ratio in each case of the total area of the disks to the area of the enclosing polygon. Which of the following is true?


 $\textbf{(A)}\ T=H=R \qquad \textbf{(B)}\ H<R=T \qquad \textbf{(C)}\ H=R<T\qquad \textbf{(D)}\ H<R<T \qquad \textbf{(E)}\ H<T<R \qquad$

                    \begin{figure}[H]
                        \centering
                        \includegraphics[width=0.5\textwidth]{Problem/AMC_10_10B_2025_12_Problem_Figure_1.png}
                    \end{figure}
                    


	\item (\textit{AMC 10 10B 2025}) The altitude to the hypotenuse of a $30{-}60{-}90^\circ$ right triangle is divided into two segments of lengths $x < y$ by the median to the shortest side of the triangle. What is the ratio $\frac{x}{x+y}$?


 $\textbf{(A)}~\frac 37 \qquad \textbf{(B)}~\frac{\sqrt3}{4} \qquad \textbf{(C)}~\frac 49 \qquad \textbf{(D)}~\frac{5}{11} \qquad \textbf{(E)}~\frac{4\sqrt3}{15}$


	\item (\textit{AMC 10 10B 2025}) Nine athletes, no two of whom are the same height, try out for the basketball team. One at a time, they draw a wristband at random, without replacement, from a bag containing $3$ blue bands, $3$ red bands, and $3$ green bands. They are divided into a blue group, a red group, and a green group. The tallest member of each group is named the group captain. What is the probability that the group captains are the three tallest athletes?


 $\textbf{(A) } \frac29 \qquad \textbf{(B) } \frac 27 \qquad \textbf{(C) } \frac9{28}\qquad \textbf{(D) } \frac13 \qquad \textbf{(E) } \frac38$


	\item (\textit{AMC 10 10B 2025}) The sum 
 \[\sum_{k=1}^{\infty} \frac{1}{k^3 + 6k^2 + 8k}\] can be expressed as $\frac{a}{b}$, where $a$ and $b$ are relatively prime positive integers. What is $a + b$?


 $\textbf{(A) } 89 \qquad \textbf{(B) } 97 \qquad \textbf{(C) } 102 \qquad \textbf{(D) } 107 \qquad \textbf{(E) } 129$


	\item (\textit{AMC 10 10B 2025}) A circle has been divided into $6$ sectors of different sizes. Then $2$ of the sectors are painted red, $2$ painted green, and $2$ painted blue so that no two neighboring sectors are painted the same color. One such coloring is shown below.


 
                    \begin{figure}[H]
                        \centering
                        \includegraphics[width=0.5\textwidth]{Problem/AMC_10_10B_2025_16_Problem_Figure_1.png}
                    \end{figure}
                    


 How many different colorings are possible?


 

$\textbf{(A)}~12\qquad\textbf{(B)}~16\qquad\textbf{(C)}~18\qquad\textbf{(D)}~24\qquad\textbf{(E)}~28$


	\item (\textit{AMC 10 10B 2025}) Consider a decreasing sequence of $n$ positive integers $x_1 > x_2 > \dotsb > x_n$ that satisfies the following two conditions:


 $\qquad\bullet$ The average (arithmetic mean) of the first $3$ terms in the sequence is $2025.$
$\qquad\bullet$ For all $4 \leq k \leq n,$ the average of the first $k$ terms in the sequence is $1$ less than the average of the first $k - 1$ terms in the sequence.


 What is the greatest possible value of ${n?}$
$\textbf{(A) } 1013 \qquad \textbf{(B) } 1014 \qquad \textbf{(C) } 1016 \qquad \textbf{(D) } 2016 \qquad \textbf{(E) } 2025$


	\item (\textit{AMC 10 10B 2025}) What is the ones digit of the sum

 \[\lfloor \sqrt{1} \rfloor + \lfloor \sqrt{2} \rfloor + \lfloor \sqrt{3} \rfloor + \dots + \lfloor \sqrt{2025} \rfloor?\](Recall that $\lfloor x \rfloor$ represents the greatest integer less than or equal to $x$.)


 $\textbf{(A) } 1\qquad\textbf{(B) } 2 \qquad\textbf{(C) }3\qquad\textbf{(D) } 5 \qquad\textbf{(E) } 8$


	\item (\textit{AMC 10 10B 2025}) A container has a $1\times 1$ square bottom, a $3\times 3$ open square top, and four congruent trapezoidal sides, as shown. Starting when the container is empty, a hose that runs water at a constant rate takes $35$ minutes to fill the container up to the midline of the trapezoids. 


 
                    \begin{figure}[H]
                        \centering
                        \includegraphics[width=0.5\textwidth]{Problem/AMC_10_10B_2025_19_Problem_Figure_1.png}
                    \end{figure}
                    


 How many more minutes will it take to fill the remainder of the container?


 $\textbf{(A) }70 \qquad \textbf{(B) }85 \qquad \textbf{(C) }90 \qquad \textbf{(D) }95 \qquad \textbf{(E) }105$


	\item (\textit{AMC 10 10B 2025}) Four congruent semicircles are inscribed in a square of side length $1$ so that their diameters are on the sides of the square, one endpoint of each diameter is at a vertex of the square, and adjacent semicircles are tangent to each other. A small circle centered at the center of the square is tangent to each of the four semicircles, as shown below.
The diameter of the small circle can be written as $(\sqrt a+b)(\sqrt c+d)$, where $a$, $b$, $c$, and $d$ are integers. What is $a + b + c + d$?


 
                    \begin{figure}[H]
                        \centering
                        \includegraphics[width=0.5\textwidth]{Problem/AMC_10_10B_2025_20_Problem_Figure_1.png}
                    \end{figure}
                    
$\textbf{(A) } 3 \qquad \textbf{(B) } 5 \qquad \textbf{(C) } 8 \qquad \textbf{(D) } 9 \qquad \textbf{(E) } 11$


	\item (\textit{AMC 10 10B 2025}) Each of the $9$ squares in a $3 \times 3$ grid is to be colored red, blue, or yellow in such a way that each red square shares an edge with at least one blue square, each blue square shares an edge with at least one yellow square, and each yellow square shares an edge with at least one red square. Colorings that can be obtained from one another by rotations and/or reflections are to be considered the same. How many different colorings are possible?


 $\textbf{(A) } 3 \qquad \textbf{(B) } 9 \qquad \textbf{(C) } 12 \qquad \textbf{(D) } 18 \qquad \textbf{(E) } 27$


	\item (\textit{AMC 10 10B 2025}) A seven-digit positive integer is chosen at random. What is the probability that the number is divisible by $11$, given that the sum of its digits is $61?$
$\textbf{(A) } \frac{3}{14} \qquad\textbf{(B) } \frac{3}{11} \qquad\textbf{(C) }\frac{2}{7}\qquad\textbf{(D) } \frac{4}{11} \qquad\textbf{(E) } \frac{3}{7}$


	\item (\textit{AMC 10 10B 2025}) A rectangular grid of squares has $141$ rows and $91$ columns. Each square has room for two numbers. Horace and Vera each fill in the grid by putting the numbers from $1$ through $141 \times 91 = 12{,}831$ into the squares. Horace fills the grid horizontally: he puts $1$ through $91$ in order from left to right into row $1$, puts $92$ through $182$ into row $2$ in order from left to right, and continues similarly through row $141$. Vera fills the grid vertically: she puts $1$ through $141$ in order from top to bottom into column $1$, then $142$ through $282$ into column $2$ in order from top to bottom, and continues similarly through column $91$. How many squares get two copies of the same number?


 $\textbf{(A) } 7 \qquad \textbf{(B) } 10 \qquad \textbf{(C) } 11 \qquad \textbf{(D) } 12 \qquad \textbf{(E) } 19$


	\item (\textit{AMC 10 10B 2025}) A frog hops along the number line according to the following rules.


 $\qquad\bullet$ It starts at $0$.


 $\qquad\bullet$ If it is at $0$, then it moves to $1$ with probability $\frac{1}{2}$ and it disappears with probability $\frac{1}{2}$.


 $\qquad\bullet$ For $n = 1, 2,$ or $3,$ if it is at $n,$ then it moves to $n+1$ with probability $\frac{1}{4},$ it moves to $n-1$ with probability $\frac{1}{4},$ and it disappears with $\qquad\phantom{\bullet}$ probability $\frac{1}{2}$.


 What is the probability that the frog reaches ${4\,?}$
$\textbf{(A) } \frac{1}{101} \qquad\textbf{(B) } \frac{1}{100} \qquad\textbf{(C) } \frac{1}{99} \qquad\textbf{(D) } \frac{1}{98} \qquad\textbf{(E) } \frac{1}{97}$


	\item (\textit{AMC 10 10B 2025}) Square $ABCD$ has sides of length $4$. Points $P$ and $Q$ lie on $\overline{AD}$ and $\overline{CD}$, respectively, with $AP=\frac{8}{5}$ and $DQ=\frac{10}{3}$. A path begins along the segment from $P$ to $Q$ and continues by reflecting against the sides of $ABCD$ (with congruent incoming and outgoing angles). If the path hits a vertex of the square, it terminates there; otherwise it continues forever. At which vertex does the path terminate?

                    \begin{figure}[H]
                        \centering
                        \includegraphics[width=0.5\textwidth]{Problem/AMC_10_10B_2025_25_Problem_Figure_1.png}
                    \end{figure}
                    
$\textbf{(A) } A \qquad \textbf{(B) } B \qquad \textbf{(C) } C \qquad \textbf{(D) } D  \qquad \textbf{(E) } \text{The path continues forever.}$


\end{enumerate}